\begin{figure}[t]
  \centering
  \fbox{
    \begin{minipage}[c][6cm][c]{0.95\columnwidth}
      \centering
      Overview
    \end{minipage}
  }
  \caption{Overview.}
  \label{fig:overview}
\end{figure}

\section{Pose Optimization and Structural Simulation} \label{sec:pose_optimization}

In this section, we present the pose optimization and structural simulation components that are responsible for refining the candidate poses generated by the initial pose sampling component to find a physically feasible configuration that closely resembles the target shape.

\subsection{Single Object Pose Optimization}
Given a candidate pose $\hat{q}_{i_{k+1}}$ for an object $\object_{i_{k+1}}$ in the inventory, the pose optimization component refines the candidate pose to find a locally optimal placement pose $q_{i_{k+1}}$ that likely leads to a physically feasible configuration and better approximates the target shape $\target$.
For a given scene $\scene_k$, a target shape $\target$, and a candidate object $\object_{i_{k+1}}$, each pose refinement can be formulated as solving a constrained optimization problem:
\begin{gather*}
    \min_{\substack{q \in \SE(3), \\ \lambda \in \real^{3n_c}}} f_\text{config}(q\mid \scene_k, \target, \object_{i_{k+1}}) + f_\text{force}(q, \lambda \mid \scene_k, \object_{i_{k+1}}) \\ 
    \text{s.t. } g_{\text{cone},j}(q,\lambda_j\mid \scene_k, \object_{i_{k+1}}) \leq 0, \quad j=1,\cdots,n_c
\end{gather*}
where $\lambda = (\lambda_1,\cdots,\lambda_{n_c}) \in \real^{3n_c}$ is the vector of contact forces at the $n_c$ contact points between the object and the scene.
Each $\lambda_j$ is represented in a 
Each $g_{\text{cone},j}(\cdot)$ is the smoothed friction cone constraints for the $j$-th contact point defined by
\begin{equation*}
    g_{\text{cone},j}(q, \lambda_j \mid \scene_k, \object_{i_{k+1}}) = \sqrt{\|J_{c,j}^t \lambda_j \|^2 + \epsilon} - \mu_j  J_{c,j}^n \lambda_j
\end{equation*}
where $\epsilon > 0$ is a smoothing parameter, $\mu_j$ is the friction coefficient, and $J_{c,j}^t$ and $J_{c,j}^n$ are the tangential and normal components of the Jacobian of the contact constraint at the $j$-th contact point.
The configuration dependent objective function $f_\text{config}(q \mid \scene_k, \target, \object_{i_{k+1}})$ is a function that evaluates how well the configuration approximates the target shape $\target$ while satisfying the physical constraints such as collision avoidance of $\object_{i_{k+1}}$ in the scene $\scene_k$.
The force dependent objective function $f_\text{force}(q,\lambda \mid \scene_k, \object_{i_{k+1}})$ is a function that evaluates whether the contact forces $\lambda$ are physically feasible and contribute to a stable configuration.
For simplicity, we will omit the dependency of $f_\text{config}$, $f_\text{force}$ and $g_{\text{cone},j}$ on $\scene_k$, $\target$, and $\object_{i_{k+1}}$ in the following subsections.

Given a fixed pose $q$, the contact forces $\lambda$ can be efficiently optimized by solving a quadratic programming problem under the second-order cone constraints.
We will describe the details of the force optimization in the next subsection.
Denoting $\lambda^*(q)$ as the optimal contact forces for a given pose $q$,  we can define the force dependent objective function as $f_\text{force}(q) := f_\text{force}(q, \lambda^*(q))$.
Moreover, denote $\nu^*(q) \in \real^{n_c}$ as the optimal dual variables corresponding to the friction cone constraints.
Then, using the \textit{Envelope theorem} \eqref{eq:envelope}, we can compute the gradient of $f_\text{force}(q)$ with respect to $q$ as follows:
\begin{equation*}
    \frac{df_\text{force}(q)}{dq} = \frac{\partial f_\text{force}(q, \lambda^*(q))}{\partial q} + \sum_{j=1}^{n_c} \nu^*_j(q) \frac{\partial g_{\text{cone},j}(q, \lambda^*_j(q))}{\partial q}.
\end{equation*}

Using the above gradient, we optimize the placement pose $q$ on SE(3) manifold based on the Riemannian trust-region method \cite{absil2007trust} combined with Riemannian BFGS Hessian approximation \cite{huang2018riemannianbfgs} to find a locally optimal placement pose $q_{i_{k+1}}$ that minimizes the objective function $f_\text{config}(q) + f_\text{force}(q)$ while satisfying the friction cone constraints $g_{\text{cone},j}(q, \lambda^*_j(q)) \leq 0$ for $j=1,\cdots,n_c$.

% In details, the configuration dependent objective function $f_\text{config}$ is defined as the following:
% \begin{equation*}
%     \begin{split}
%         &f_\text{config}(q \mid \scene_k, \target, \object_{i_{k+1}}) = \\
%         &\quad f_\text{height}(q) + f_\text{inclusion}(q \mid \target, \object_{i_{k+1}}) 
%         + f_\text{collision}(q \mid \scene_k, \object_{i_{k+1}}).
%     \end{split}
% \end{equation*}
% where the three terms are functions that evaluate the following aspects of the configuration:
% \begin{itemize}
%   \item $f_\text{height}(q)$: the height of the object $\object_{i_{k+1}}$ at the pose $q$ to encourage the object to be placed at a lower height to improve stability,
%   \item $f_\text{inclusion}(q \mid \target, \object_{i_{k+1}})$: the Hausdorff distance between the current configuration and the target shape $\target$ to encourage the configuration to better approximate the target shape, and
%   \item $f_\text{collision}(q \mid \scene_k, \object_{i_{k+1}})$: penalizes configurations that result in collisions between the object and the scene.
% \end{itemize}

\subsection{Force Optimization}
For indices $l, m \in \placed_k$, let $I_{l,m} \subset \{1,\cdots,n_c\}$ be the set of contact point indices associated with the objects $\object_l$, $\object_m$, and let $I_{l} := \bigcup_{m \in \placed_k} I_{l,m}$ be the set of contact point indices associated with the object $\object_l$.
Then, the force dependent objective $f_\text{force}(q,\lambda)$ is defined as the following:
\begin{equation}
\begin{split}
  f_\text{force}(q,\lambda) 
   =& \sum_{l \in \placed_k} \|J_{c,I_l}^\top(q)  \lambda_{I_l} + \lambda_{\text{ext},l}\|^2 + \rho\sum_{j = 1}^{n_c} \|\lambda_j\|^2 \\
   &+ \sum_{l \in \placed_k}\sum_{j \in I_{l,i_{k+1}}} \sd_j(q)^2 \|\lambda_j\|^2
\end{split}
\end{equation}
where $J_{c,I_l}(q) \in \real^{3|I_l| \times 6}$ is the Jacobian of the contact constraints associated with the object $\object_l$ at the contact points in $I_l$, $\lambda_{I_l} \in \real^{3|I_l|}$ is the vector of contact forces at the contact points in $I_l$, and $\lambda_{\text{ext},l} \in \real^6$ is the external wrench acting on the object $\object_l$.
The first term in $f_\text{force}(q,\lambda)$ penalizes the violation of the static equilibrium constraints for each object in the scene.
The second term with $\rho > 0$ is a regularization term to avoid unnecessarily large contact forces which is not physically meaningful.
The third term penalizes the complementarity conditions on the contact forces between the candidate object $\object_{i_{k+1}}$ and the objects in the scene $\placed_k$ based on the signed distance $\sd_j(q)$ between the two objects at the $j$-th contact point which can be computed as described in Section \ref{sec:preliminaries}.
Note that $\sd_j(q)$ is differentiable which allows us to apply \textit{Envelope theorem} to compute the gradient of $f_\text{force}(q)$ with respect to $q$ as described in the previous subsection.



\subsection{Structural Simulation}
